\begin{abstract}
Choosing how to spend an annotation budget when deploying an object detector is hard: an open-vocabulary detector can be used zero-shot, used as a teacher whose pseudo-labels train a small student, or replaced by a detector fine-tuned on manual labels. We present a budget-controlled comparison of these strategies on two public domains, HomeObjects-3K (common household nouns) and Construction-PPE (domain-specific classes, including absence of equipment), with nested labeled subsets of 0--400 images, matched training compute and a single evaluator. The best strategy depends on the domain: on HomeObjects, zero-shot YOLOE-11s (55.3 mAP$_{50}$) was not surpassed by any trained YOLO11n student up to 400 labels (28.4), although a compute ablation shows the gap narrows from 27 to about 14 points with more training; on Construction-PPE, zero-shot reaches only 23.6 and fine-tuning on 10 labeled images already exceeds it (27.4), reaching 41.1 at 400. Pseudo-label distillation at best matched the teacher, and mixing noisy pseudo-labels with 200 gold labels lowered accuracy by 9.7 points compared with using the gold labels alone. We further propose a calibration of the pseudo-label threshold on the labeled images (no measurable effect here), a learning-curve-based strategy selector (correct in all 12 low-difficulty test cases, which contain no real crossover), and quantify that 50 labeled images estimate zero-shot mAP$_{50}$ to within about 2.5--3.4 points. Results are small-scale (two datasets, two seeds, CPU-limited students) and should be validated with stronger teachers and larger budgets.
\end{abstract}

\begin{IEEEkeywords}
open-vocabulary object detection, annotation budget, pseudo-labeling, teacher--student learning, fine-tuning, strategy selection
\end{IEEEkeywords}
